% Page 065 — page imprimée 62
% Anecdotes, Témoignages, Souvenirs…


\sectiondate{Anecdotes, Témoignages, Souvenirs…}{}


\medskip

\noindent{\bfseries EVA VERDIER} mère de Madeleine qui avait épousé Emilien Martin (originaire de
Rousses à 3km du Villaret dans la vallée du Bétuzon) raconte, à 104 ans, en 1986, à la
Bragousette où elle passa ses dernières années chez sa fille :

\medskip

\subsection{La soie}

\begin{quote}
\textit{« Quand j’étais jeune j’habitais aux Salides, un petit hameau sur la pente sous le col Salidès.
Les mûriers étaient nombreux dans la région et on élevait des vers à soie dans beaucoup de
fermes. Les cocons étaient vendus à St. André de Valborgne, et je faisais éclore quelques
onces de « graine ». Dès que j’avais 5kg de cocons je les portais à la filature à Saint André. »}
\end{quote}

\begin{quote}
\textit{« La Fabrique » à Meyrueis était alors une grande magnanerie, puis elle fut transformée en
filature. »}
\end{quote}

\medskip

\subsection{Un assassinat en 1826}

Celui de Monsieur de Fenouillet, Baron des Fons, (Les Fons est une grosse propriété à
quelques km. de Cabrillac sur la route de Saint André de Valborgne). C’est une histoire que
lui avait racontée son grand père :

\begin{quote}
\textit{« Monsieur de Fenouillet, propriétaire terrien vivant du revenu de ses terres cultivées par
des fermiers faisait confectionner par sa lingère habitant aux Salides (hameau en contrebas
du col Salidès) des sacs d’étoffe dans les quels il serrait les rouleaux d’écus qu’il amassait. Il
les cachait dans une pièce spéciale du château des Fons (était-ce son bureau ou sa cave, je ne
sais).}
\end{quote}

\begin{quote}
\textit{Trois jeunes gens des Salides décidèrent de le dévaliser. L’un d’eux, courtisant la lingère,
manœuvre habilement pour localiser le trésor, lui demandant d’abord de soustraire elle
même quelques sachets et se heurtant à un refus finit, après lui avoir demandé quel est son
travail principal, à quoi étaient destinés les sacs qu’elle cousait, comment et où ils étaient
rangés, à obtenir une nuit un rendez vous amoureux. Il la prie de laisser le soir une porte du
château ouverte.}
\end{quote}

\begin{quote}
\textit{Il pénétra alors la nuit avec ses deux acolytes et ils arrivèrent à la cachette, lorsque le baron
réveillé en sursaut par un bruit insolite se leva et leur barra la route. Pris de panique les trois
voleurs lui tranchèrent la gorge. Ils furent rapidement identifiés et condamnés à être
décapités.}
\end{quote}

\begin{quote}
\textit{La guillotine fut installée au Pompidou. Un jeune berger des Salides décida d’assister à
l’exécution pour « voir tomber les têtes ». Ce fut si horrible qu’il en perdit la raison. »}
\end{quote}

% Page 066 — page imprimée 63
% Suite de « Anecdotes, Témoignages, Souvenirs… »

Voilà donc ce récit conservé dans une famille, par la tradition orale. En fait il semble que si
cet assassinat a bien eu lieu, les pièces officielles, conservées à l’occasion du procès des
assassins les détails furent tout autres.

M. Maurice Abric de Fernouillet fut bien assassiné par trois hommes, le 16 janvier 1826 mais
au Château de l’Hom, une autre de ses propriétés, au Follaquier, hameau près de la route de
Saint André de Valborgne à quelques km des Salides et non pas aux Fons.
Sans doute y eut-il vol mais M. de Fenouillet, noble de fraîche date faisait preuve
d’arrogance, humiliant ses fermiers, corrigeant lui-même à coups de bâton deux enfants de 10
ans ! fils de son fermier, qui avaient dérobé quelques une de ses poires… n’hésitant pas à
couper les oreilles de chèvres qui avaient mangé ses châtaignes ou écorcé de jeunes
châtaigniers :

\begin{quote}
\textit{« …ce type là, raconte un habitant de Bassurel dans les années 1980, il aimait pas les gens.
Et alors les gens faisaient sortir leurs chèvres, comme ça. Et quand il trouvait des chèvres, il
leur coupait les oreilles ! Et quand ils ont vu ça, un jour ils l’ont tué, les voisins. Et je crois
qu’il est enterré là bas à l’Hom… Oh c’était une sale bête ! Il pouvait pas voir ses voisins,
quoi. »}
\end{quote}

Les trois assassins furent bel et bien condamnés à mort et guillotinés sur la place du
Pompidou le 7 mars 1827.

Cette triste histoire a tout de même causé la mort de quatre personnes, semé la désolation dans
plusieurs familles et laissé une dizaine d’orphelins dont certains n’ont pas connu leur père.

\begin{flushright}
Almanach de la Vallée Borgne 2006
\end{flushright}
